% TEMPLATE-MILC-v3.tex - plantilla maestra UNICA de las guias MILC v3.
%
% La usan las cuatro familias de guias, cada una con su driver:
%   · Grados 6-11          -> scripts/build-guias-g11.py
%   · Bebras               -> scripts/build-guias-bebras.py
%   · Semillero            -> scripts/build-guias-semillero.py
%   · Territorio interior  -> scripts/build-guias-territorio-interior.py
%
% Antes habia tres copias casi identicas (~400 lineas iguales, 6 distintas):
% cada mejora habia que aplicarla tres veces, y en la practica no se hacia
% (el motor de recursos vivia solo en dos de ellas). Lo que varia entre
% familias esta parametrizado en 6 placeholders que cada driver llena
% (se nombran aqui SIN su sintaxis real: el builder reemplaza por texto plano,
% tambien dentro de los comentarios, y un valor multilinea rompe el preambulo):
%
%   HEADER_IZQ        encabezado izquierdo de pagina
%   HEADER_DER        encabezado derecho de pagina
%   PORTADA_ETIQUETA  rotulo sobre el numero grande (GRADO/MOMENTO/MODULO)
%   PORTADA_SERIE     linea de serie de la portada
%   PORTADA_ID        identificador de la guia en la portada
%   PORTADA_CIRCULO   el circulito de grado (°), o vacio si no aplica
%   PDF_TITULO        titulo en texto plano (sin \textbf ni comillas TeX) para
%                     los metadatos del PDF (/Title); lo produce lib_xelatex.texto_pdf
%
% Composicion (auditoria 2026-09-03): las bandas de estacion piden espacio
% (\Needspace*) y prohiben el salto justo despues (\nopagebreak), asi no quedan
% huerfanas al pie; las cajas largas (softbox, drawbox, infoband, puentebox) son
% `breakable`, asi una caja de media pagina ya no arrastra todo a la siguiente
% dejando la anterior al 40 %. Todo texto sobre mostaza o verde va en milcNegro
% (blanco sobre #E5B400 da 1,93:1 y sobre #58A923 2,95:1; ninguno pasa WCAG AA).
% El resto de claves las documenta content/guias/_SCHEMA.md. El script valida
% que no queden placeholders sin reemplazar antes de escribir el archivo final.

% Etiquetado PDF (latex-lab): /Lang, estructura y texto alternativo de las
% figuras, para lectores de pantalla. Debe ir ANTES de \documentclass.
\DocumentMetadata{lang=es-CO,tagging=on}
\documentclass[11pt,letterpaper]{article}
% headheight=24pt: la cabecera derecha lleva dos lineas (identidad de la guia
% y estacion actual). headsep se acorta en la misma medida para que el bloque
% de texto quede exactamente donde estaba con headheight=12pt.
\usepackage[margin=2.0cm,headheight=24pt,headsep=13pt]{geometry}
\usepackage[table]{xcolor}
\usepackage{fontspec}
\usepackage[spanish]{babel}
\usepackage{tikz}
% Librerías TikZ para los diagramas de las guías (bloque `recursos:`):
%  · babel  → babel en español vuelve activos caracteres como «>», y TikZ los usa
%             en sus opciones (>=stealth, ->). Sin esta librería, cualquier
%             diagrama con flechas revienta con «\language@active@arg> has an
%             extra }». Debe ir ANTES de que se dibuje nada.
%  · positioning        → sintaxis «below left=2pt and 3pt».
%  · decorations.pathreplacing → llaves ({brace}) para señalar rangos.
\usetikzlibrary{babel,positioning,decorations.pathreplacing}
\usepackage{graphicx}
% Circuitos: el trabajo de electrónica se hace hoy con micro:bit y sus sensores
% integrados (MakeCode, sin cableado), así que no se necesitan esquemáticos.
% Si algún día entran montajes con protoboard, basta descomentar esta línea:
% los diagramas .tex/.tikz declarados en `recursos:` ya se incrustan solos.
% \usepackage{circuitikz}
\usepackage[most]{tcolorbox}
\usepackage{tabularx}
\usepackage{array}
\usepackage{fancyhdr}
\usepackage{titlesec}
\usepackage[document]{ragged2e}  % bandera a la derecha en todo el cuerpo
\usepackage{enumitem}
\usepackage{microtype}
\usepackage{etoolbox}
\usepackage{needspace}
% hyperref al final del preambulo: metadatos del PDF (titulo, autor, idioma,
% familia), marcadores por estacion (\pdfbookmark) y enlaces sin recuadro.
\usepackage[hidelinks,
  pdftitle={El cuerpo comparado: quién fabricó la vara con que te mides},
  pdfauthor={PhD. Álvaro Cárdenas Orozco},
  pdfsubject={Territorio interior · Educación socioemocional · Grado 8° · Momento 7 · MILC v3},
  pdflang={es-CO},
  bookmarksopen=true,bookmarksnumbered=false]{hyperref}

% Helvetica Neue no trae la flecha U+2192, asi que cada «→» escrito literalmente
% en el YAML se compilaba a NADA, en silencio (462 flechas en 61 guias: rutas del
% tipo «paleta → tipografia → lockup» salian rotas). Mapeamos el caracter a su
% version matematica, que si tiene glifo. Asi el autor puede seguir escribiendo
% «→» comodamente en el YAML.
\usepackage{newunicodechar}
\newunicodechar{→}{\ensuremath{\rightarrow}}
% Mismo problema con los simbolos de marca y vinetas: el «✓» de «una palomita ✓»
% salia como cuadrito vacio en el PDF de todas las guias que lo usaban.
\usepackage{pifont}
\newunicodechar{✓}{\ding{51}}
\newunicodechar{✗}{\ding{55}}
\newunicodechar{✔}{\ding{52}}
\newunicodechar{•}{\textbullet}
\newunicodechar{≥}{\ensuremath{\geq}}
\newunicodechar{≤}{\ensuremath{\leq}}

\setmainfont{Helvetica Neue}
\setsansfont{Helvetica Neue}
\setlength{\parindent}{0pt}
\setlength{\parskip}{6pt}
% Sin particion de palabras: en 6.º-7.º una palabra cortada por guion al final
% de linea («Comuni-cacion») frena la lectura mucho mas que una linea corta.
\hyphenpenalty=10000
\exhyphenpenalty=10000
\pretolerance=2000
\tolerance=3000
\setlength{\emergencystretch}{3em}
\linespread{1.10}
\renewcommand{\arraystretch}{1.25}
\setlist{leftmargin=5mm,itemsep=1mm,topsep=1mm}
\newcolumntype{Y}{>{\RaggedRight\arraybackslash}X}

\definecolor{milcMagenta}{HTML}{E6007E}
\definecolor{milcVino}{HTML}{5A0038}
\definecolor{milcTurquesa}{HTML}{0093A5}
\definecolor{milcVerde}{HTML}{58A923}
\definecolor{milcMostaza}{HTML}{E5B400}
\definecolor{milcOcre}{HTML}{B86600}
\definecolor{milcNegro}{HTML}{241816}
\definecolor{milcGris}{HTML}{F7F7F4}
\definecolor{milcLinea}{HTML}{E8E2DD}
\definecolor{milcGradePrimary}{HTML}{0D9488}
\definecolor{milcGradeDark}{HTML}{115E59}
\definecolor{milcGradeSoft}{HTML}{CCFBF1}
\definecolor{milcGradeAccent}{HTML}{CFE7FF}
\definecolor{milcGradeText}{HTML}{FFFFFF}
\color{milcNegro}

\pagestyle{fancy}
\fancyhf{}
% Cabecera de dos lineas: identidad de la guia arriba y, a la derecha y en
% gris, la estacion en curso (\markright desde \milcestacion). Antes de la
% primera estacion la segunda linea va vacia; el \strut mantiene la altura
% para que la linea de arriba no baile entre paginas.
\lhead{\small Territorio interior · Educación socioemocional\\[-1pt]{\footnotesize\strut}}
\rhead{\small Grado 8° · Momento 7 · MILC v3\\[-1pt]{\footnotesize\strut\color{milcNegro!65}\rightmark}}
\cfoot{\small\thepage}
\renewcommand{\headrulewidth}{0pt}

% Sobre mostaza (#E5B400) y verde (#58A923) el texto va en milcNegro; sobre el
% resto de la paleta, en blanco. Se usa dentro de fonttitle/fontupper, donde un
% \color posterior manda sobre coltitle.
\newcommand{\milctintasobre}[1]{%
  \ifboolexpr{test{\ifstrequal{#1}{milcMostaza}} or test{\ifstrequal{#1}{milcVerde}}}%
    {\color{milcNegro}}{\color{white}}}

\newtcolorbox{titlebox}[1]{
  enhanced,colback=#1,colframe=#1,arc=7mm,boxrule=0pt,
  left=7mm,right=7mm,top=3mm,bottom=3mm,
  before skip=8pt,after skip=8pt,
  fontupper=\sffamily\bfseries\Large\color{white}
}

\newtcolorbox{softbox}[3]{
  enhanced,breakable,pad at break=3mm,
  colback=#2,colframe=#1,borderline west={4pt}{0pt}{#1},
  arc=2mm,boxrule=.4pt,left=4mm,right=4mm,top=3mm,bottom=3mm,
  before skip=6pt,after skip=8pt,title={#3},coltitle=white,
  colbacktitle=#1,fonttitle=\sffamily\bfseries\milctintasobre{#1},fontupper=\RaggedRight,
  attach boxed title to top left={xshift=3mm,yshift=-2mm},
  boxed title style={arc=2mm, boxrule=0pt}
}

\newtcolorbox{drawbox}[2]{
  enhanced,breakable,pad at break=4mm,
  colback=white,colframe=#1,arc=4mm,boxrule=1pt,
  left=5mm,right=5mm,top=4mm,bottom=4mm,before skip=8pt,after skip=8pt,
  title={#2},coltitle=white,colbacktitle=#1,fonttitle=\sffamily\bfseries\milctintasobre{#1},
  fontupper=\RaggedRight,
  attach boxed title to top left={xshift=4mm,yshift=-2mm},
  boxed title style={arc=3mm, boxrule=0pt}
}

\newtcolorbox{infoband}[1]{
  enhanced,breakable,pad at break=2mm,colback=#1!8,colframe=#1!50,arc=2mm,boxrule=.5pt,
  left=4mm,right=4mm,top=2mm,bottom=2mm,before skip=4pt,after skip=4pt,
  fontupper=\small\RaggedRight
}

% Puente narrativo entre secciones (contrato editorial Conéctate).
% Da continuidad: "Acabas de X. Ahora vas a Y porque Z."
% Visualmente: banda izquierda en color del grado, texto en itálica pequeña.
\newtcolorbox{puentebox}{
  enhanced,breakable,pad at break=2mm,colback=milcGradeSoft,colframe=milcGradeDark,
  borderline west={3pt}{0pt}{milcGradeDark},
  arc=0mm,boxrule=0pt,left=5mm,right=4mm,top=2.5mm,bottom=2.5mm,
  before skip=4pt,after skip=8pt,
  fontupper=\itshape\small\color{milcNegro}\RaggedRight
}
% La flecha va en modo matematico a proposito: Helvetica Neue NO tiene el glifo
% U+2192, asi que \textrightarrow se compilaba a NADA («Missing character: There
% is no → in font Helvetica Neue») y el puente salia sin flecha. En modo
% matematico la flecha viene de la fuente matematica y si se dibuja.
\newcommand{\puente}[1]{%
  \begin{puentebox}%
    \textcolor{milcGradeDark}{$\rightarrow$}\hspace{2mm}#1%
  \end{puentebox}%
}

\newcommand{\linead}[1]{\noindent\color{milcLinea}\rule{#1}{0.5pt}\color{milcNegro}}
\newcommand{\respuesta}[1]{\par\vspace{2mm}\foreach \n in {1,...,#1}{\noindent\color{milcLinea}\rule{\linewidth}{0.5pt}\par\vspace{5mm}}\color{milcNegro}}
\newcommand{\checkbox}{\textcolor{milcGradeDark}{\fbox{\phantom{x}}}\hspace{5pt}}

% ─── Listas aireadas para las guías socioemocionales ────────────────────
% Componer las verificaciones como «\checkbox texto\\» deja el texto que salta
% de línea debajo del cuadrito y apelmaza el bloque. Estos entornos alinean la
% continuación y airean los ítems. Los usa Territorio Interior; cualquier
% familia puede hacerlo.
\newlist{checklista}{itemize}{1}
\setlist[checklista]{
  label=\textcolor{milcGradeDark}{\fbox{\phantom{x}}},
  leftmargin=9mm, labelsep=3.2mm, itemsep=2.4mm,
  topsep=2.5mm, parsep=0pt, partopsep=0pt, before=\RaggedRight
}
\newlist{pasoslista}{enumerate}{1}
\setlist[pasoslista]{
  label=\textcolor{milcGradeDark}{\sffamily\bfseries\arabic*.},
  leftmargin=8mm, labelsep=3mm, itemsep=2.8mm,
  topsep=1mm, parsep=0pt, partopsep=0pt, before=\RaggedRight
}

\newcommand{\phasecard}[4]{
\begin{tcolorbox}[enhanced,colback=#3,colframe=#2,arc=3mm,boxrule=.5pt,height=3.05cm,valign=center,halign=center,left=1.5mm,right=1.5mm,top=2mm,bottom=2mm]
{\sffamily\bfseries\fontsize{22}{24}\selectfont\textcolor{#2}{#1}}\par
{\sffamily\bfseries\small #4}
\end{tcolorbox}
}

% ── Piezas del rediseno 2026 (legibilidad 6.º-11.º) ───────────────────
% Banda de estacion: reemplaza el titlebox macizo. Titulo a la izquierda,
% dato practico (tiempo/modalidad) a la derecha, en una sola linea.
% Nunca huerfana: pide 9 lineas libres (o salta de pagina rellenando con
% \vfill) y prohibe el salto justo despues de la banda. Nueve y no seis:
% la caja partible que sigue necesita ~80pt (titulo adosado, relleno y dos
% lineas) para arrancar en la misma pagina; con menos, tcolorbox la manda
% entera a la siguiente y la banda se queda sola. El penalty 10000 va
% en `after`, ANTES del espacio posterior: un `after skip` (aunque sea 0pt)
% es un \vskip tras la caja, y ese glue es un punto de corte legal que un
% \nopagebreak escrito despues de \end{tcolorbox} ya no cierra. Cada banda
% deja marcador PDF y actualiza la cabecera derecha.
\newcounter{milcestacion}
\newcommand{\milcestacion}[2]{%
\Needspace*{9\baselineskip}%
\stepcounter{milcestacion}%
\pdfbookmark[1]{#2}{milcestacion\themilcestacion}%
\markright{#2}%
\begin{tcolorbox}[enhanced,colback=#1,colframe=#1,arc=2mm,boxrule=0pt,
  left=5mm,right=5mm,top=2.6mm,bottom=2.6mm,before skip=11pt,
  after={\par\nopagebreak\vskip7pt}]
{\sffamily\bfseries\fontsize{15}{18}\selectfont\milctintasobre{#1}#2}
\end{tcolorbox}}

% Tarjeta de paso numerada: sustituye las tablas «Pilar / Aplicacion», que
% eran formato de consulta docente, no de estudiante. El numero va en un
% circulo sobre el borde para que la secuencia se lea de un vistazo.
\newcommand{\milcpaso}[3]{%
\Needspace{4\baselineskip}%
\begin{tcolorbox}[enhanced,colback=white,colframe=milcLinea,arc=1.5mm,boxrule=.9pt,
  left=14mm,right=4mm,top=3mm,bottom=3mm,before skip=4pt,after skip=4pt,
  fontupper=\RaggedRight,
  overlay={\node[anchor=north west,circle,fill=#2,text=white,inner sep=0pt,
    minimum size=7.4mm,font=\sffamily\bfseries\small]
    at ([xshift=3.6mm,yshift=-3.2mm]frame.north west) {#1};}]
#3
\end{tcolorbox}}

% Casilla marcable de tamano util con lapiz (la anterior era \fbox{\phantom{x}},
% de alto variable segun el texto que la seguia).
\newcommand{\milcmarca}{\framebox[3.8mm]{\rule{0pt}{3.4mm}}}

% Fila marcable de lista suelta (checks de la estacion 1).
\newcommand{\milccheckrow}[1]{%
\par\vspace{1.8mm}\noindent
\makebox[6.4mm][l]{\raisebox{-.55mm}{\textcolor{milcNegro}{\milcmarca}}}%
\parbox[t]{\dimexpr\linewidth-6.4mm\relax}{\RaggedRight #1}\par}

% Celda marcable dentro de la rubrica (tabla de autoevaluacion). Misma
% sangria francesa, pero pensada para vivir dentro de una columna X.
\newcommand{\milcopcion}[1]{%
\makebox[5.2mm][l]{\raisebox{-.55mm}{\textcolor{milcNegro}{\milcmarca}}}%
\parbox[t]{\dimexpr\linewidth-5.2mm\relax}{\RaggedRight #1}}

% ── Recursos visuales de la guía ──────────────────────────────────────
% Declarados en el bloque `recursos:` del YAML y emitidos por el builder.
% \guiaFigura   → imagen rasterizada o PDF (foto, carta celeste, montaje).
% \guiaDiagrama → fuente TikZ/circuitikz (.tex) incrustada como vector:
%                 es la vía para circuitos y esquemáticos editables.
% [#1] = texto alternativo (alt) · #2 = ruta relativa al .tex · #3 = pie
% El alt llega al PDF etiquetado (/Alt) y lo lee el lector de pantalla.
\newcommand{\guiaFigura}[3][]{%
\begin{center}
\includegraphics[width=0.88\linewidth,height=0.38\textheight,keepaspectratio,alt={#1}]{#2}\\[1.5mm]
{\footnotesize\itshape\textcolor{milcNegro!75}{#3}}
\end{center}
}

\newcommand{\guiaDiagrama}[2]{%
\begin{center}
\input{#1}\\[1.5mm]
{\footnotesize\itshape\textcolor{milcNegro!75}{#2}}
\end{center}
}

\begin{document}
\thispagestyle{empty}

% ── Portada ───────────────────────────────────────────────────────────
% Antes: un rectangulo de color a pagina completa. Se veia bien en pantalla,
% pero en la fotocopiadora del colegio se comia el toner y salia gris sucio.
% Ahora la banda ocupa el tercio superior y el resto va en blanco.
\begin{tikzpicture}[remember picture,overlay]
  \path[fill=milcGradePrimary]
    (current page.north west) rectangle ([yshift=-10.6cm]current page.north east);

  \node[anchor=north west,text=milcGradeText] at ([xshift=2.0cm,yshift=-1.55cm]current page.north west)
    {\sffamily\bfseries\fontsize{12}{14}\selectfont MOMENTO};
  \node[anchor=north west,text=milcGradeText,opacity=.85] at ([xshift=2.0cm,yshift=-2.35cm]current page.north west)
    {\sffamily\bfseries\fontsize{8.5}{10}\selectfont TERRITORIO INTERIOR · SOCIOEMOCIONAL};
  \node[anchor=north east,text=milcGradeText,opacity=.85,align=right] at ([xshift=-2.0cm,yshift=-1.55cm]current page.north east)
    {\sffamily\bfseries\fontsize{9}{12}\selectfont I.E. Sor María Juliana\\Cartago, Valle del Cauca};

  \node[anchor=west,text=milcGradeText] (gradonum) at ([xshift=1.85cm,yshift=-7.1cm]current page.north west)
    {\sffamily\bfseries\fontsize{112}{112}\selectfont 7};
  % El circulito de grado (°) solo aplica a las guias de grado («11°»); en
  % Bebras y Semillero el numero grande es un momento/modulo, no un grado.
  % Se posiciona relativo al nodo `gradonum`, no en coordenadas absolutas.
  

  \node[anchor=west,text=milcGradeText,text width=11.4cm,align=left] at ([xshift=1.5cm,yshift=0.5cm]gradonum.east)
    {
     {\sffamily\bfseries\fontsize{9}{11}\selectfont\MakeUppercase{Territorio interior · Socioemocional}}\\[3mm]
     {\sffamily\bfseries\fontsize{21}{25}\selectfont\setlength{\baselineskip}{27pt}El cuerpo\\[4pt]comparado\\[4pt]quién fabricó la vara}\\[3.5mm]
     {\sffamily\bfseries\fontsize{15}{18}\selectfont Grado 8° · Momento 7}
    };
\end{tikzpicture}

\vspace*{9.4cm}
\vfill

{\sffamily\bfseries\fontsize{8}{10}\selectfont\textcolor{milcGradeDark}{LO QUE TE LLEVAS HOY}}

\begin{tcolorbox}[enhanced,colback=white,colframe=milcNegro,arc=2mm,boxrule=1.6pt,
  left=5mm,right=5mm,top=4mm,bottom=4mm,before skip=3pt,after skip=12pt,
  fontupper=\RaggedRight\fontsize{11}{15.5}\selectfont]
Tu línea de la vara: tres épocas con el cuerpo que se consideraba correcto en cada una y quién lo decidía, tu propia lista de comparaciones de una semana con su fuente, y las tres reglas que vas a aplicarle a tu manera de mirar y de hablar de los cuerpos.
\end{tcolorbox}

\begin{tcolorbox}[enhanced,colback=milcGris,colframe=milcGris,arc=2mm,boxrule=0pt,
  left=5mm,right=5mm,top=3.5mm,bottom=3.5mm,after skip=12pt]
{\sffamily\bfseries\fontsize{8}{10}\selectfont\textcolor{milcNegro!55}{ESTA GUÍA ES DE}}\\[3mm]
\begin{tabularx}{\linewidth}{X p{3.2cm} p{3.2cm}}
\linead{\linewidth} & \linead{3.0cm} & \linead{3.0cm}\\[-1mm]
{\fontsize{8.5}{10}\selectfont\textcolor{milcNegro!55}{Nombre completo}} &
{\fontsize{8.5}{10}\selectfont\textcolor{milcNegro!55}{Grupo}} &
{\fontsize{8.5}{10}\selectfont\textcolor{milcNegro!55}{Fecha}}\\
\end{tabularx}
\end{tcolorbox}

\vfill

\noindent\textcolor{milcLinea}{\rule{\linewidth}{1pt}}\\[2mm]
{\sffamily\bfseries\fontsize{9.5}{12}\selectfont Autor: Dr. Álvaro Cárdenas Orozco}\\[1mm]
{\sffamily\fontsize{8.5}{11}\selectfont\textcolor{milcNegro!55}{Metodología MILC · Apertura ancestral · Imagen corporal · Triángulo Dussel-Estoicismo-Floridi}}

\newpage

\section*{Apertura: El cuerpo comparado: quién fabricó la vara con que te mides}

\begin{softbox}{milcGradeDark}{milcGradeSoft}{Saber ancestral}
Colombia \textbf{aprendió} a mirar los cuerpos, y se puede rastrear dónde. Durante buena parte del siglo XX, la revista \textbf{Cromos} ---de circulación nacional--- \textbf{promovió ideales estéticos sobre el cuerpo}, ligando moda, buen gusto, elegancia y feminidad. Hay estudios que analizan sus páginas entre \textbf{1916 y 1929} y muestran cómo la revista fue transformando las imágenes de lo que se consideraba bello. \textbf{Y hay un segundo aparato, todavía más potente: el Reinado Nacional de Belleza}, cuyos discursos en Cromos entre \textbf{1957 y 1962} han sido estudiados como la construcción del \textbf{mito del 90-60-90} ---tres números presentados como si fueran una medida de la belleza---. \textbf{Fíjate en lo que eso significa.} \emph{Esos números no salieron de la naturaleza:} \textbf{los publicó una revista, en unos años concretos, y antes de eso nadie en Colombia se medía así}. \textbf{Y hay algo más que muestran esos estudios:} los cambios en el ideal responden también a \textbf{intereses económicos} ---el crecimiento de las industrias textil, cosmética y de la higiene empujó una idea de belleza que se pudiera comprar---. \emph{Es decir: la vara con que hoy se mide un cuerpo tiene autor, tiene fecha y tiene quién gana con ella.}
\end{softbox}

\begin{softbox}{milcTurquesa}{milcGris}{Saber contemporáneo}
¿Y por qué nos comparamos, si sabemos que hace daño? \textbf{Porque no tenemos otra manera de saber cómo estamos.} En 1954, \textbf{Leon Festinger} propuso la \textbf{teoría de la comparación social} y arrancó con una idea simple: en el ser humano hay \textbf{un impulso a evaluar sus opiniones y sus capacidades}. \textbf{Y su segunda hipótesis es la que nos importa hoy:} cuando \textbf{no existe una medida objetiva} para evaluarse, \textbf{la gente se evalúa comparándose con otros} (Festinger, 1954). \emph{Piénsalo con un ejemplo:} para saber si corres rápido hay un cronómetro ---una medida objetiva---. \textbf{Para saber si tu cuerpo «está bien» no hay ningún cronómetro}, así que el único instrumento disponible son \textbf{los demás}. \textbf{Y Festinger añadió algo decisivo:} \emph{la gente tiende a comparse menos a medida que la diferencia con el otro se hace más grande}. \textbf{Es decir, tenemos un freno natural: no nos comparamos con quien está lejísimos de nosotros.} \emph{Ahora piensa en lo que hace una red social: te pone al lado, todos los días, a personas con las que jamás te habrías comparado ---y muchas veces ni siquiera son personas: son fotos escogidas entre cien y arregladas---.} \textbf{La red no inventó la comparación: le quitó el freno.}
\end{softbox}

\begin{softbox}{milcMostaza}{milcGris}{Pregunta-puente}
\textit{El 90-60-90 lo publicó \textbf{una revista, en unos años concretos}, y antes nadie se medía así. \textbf{¿Con quién te comparaste ayer} \ldots \textbf{y quién decidió que esa fuera la vara}?}
\end{softbox}

\begin{softbox}{milcVerde}{milcGris}{Saber y saber hacer}
Al terminar podrás: (1) \textbf{mostrar}, con una línea de tiempo, que el cuerpo «correcto» cambia según la época y que alguien lo decide; (2) \textbf{explicar} por qué comparse es normal ---y por qué las redes le quitan el freno que Festinger describió---; (3) \textbf{registrar} tus propias comparaciones durante una semana e identificar de dónde salen; (4) \textbf{aplicar} tres reglas concretas a tu manera de mirar y de hablar de los cuerpos, empezando por la tuya.
\end{softbox}



\small
\begin{tabularx}{\textwidth}{>{\bfseries}p{3.0cm}Y}
\rowcolor{milcGradePrimary!12}Autor & Dr. Álvaro Cárdenas Orozco\\
DBA & Comprende los fenómenos que afectan a su territorio, regula sus emociones ante ellos y acompaña a otros con empatía (transversal 6°--11°, articulado con la política «Escuela, territorio de vida» del MEN, Resolución 006519 de 2025, y la Ley 1620 de 2013)\\
Referentes & Primera ayuda psicológica (OMS, 2011/2012) · Directrices IASC (2007) · USGS y Servicio Geológico Colombiano · Pueblo nasa y la minga (CRIC) · Dussel · Estoicismo · Floridi\\
Producto final & Tu línea de la vara: tres épocas con el cuerpo que se consideraba correcto en cada una y quién lo decidía, tu propia lista de comparaciones de una semana con su fuente, y las tres reglas que vas a aplicarle a tu manera de mirar y de hablar de los cuerpos.\\
\end{tabularx}
\normalsize

\puente{En el momento 6 viste lo que se compra para pertenecer. \textbf{Hoy toca lo que no se puede comprar ni devolver:} el cuerpo con el que a uno le tocó.}

\milcestacion{milcGradeDark}{Ruta MILC: así vamos a aprender}
\begin{tabularx}{\textwidth}{>{\centering\arraybackslash}X>{\centering\arraybackslash}X>{\centering\arraybackslash}X>{\centering\arraybackslash}X}
\phasecard{1}{milcVerde}{milcGris}{Escucha\\[-1mm]\small Observo, escucho y pregunto.} &
\phasecard{2}{milcTurquesa}{milcGris}{Sistematización\\[-1mm]\small Miro la vara, no el cuerpo.} &
\phasecard{3}{milcMagenta}{milcGris}{Praxis\\[-1mm]\small Construyo y aplico.} &
\phasecard{4}{milcVino}{milcGris}{Evaluación\\[-1mm]\small Reflexión liberadora.}
\end{tabularx}


\puente{Primero, lo que se dice de los cuerpos en tu salón ---sin nombres--- y lo que te dices a ti mismo, que suele ser lo más duro.}

\milcestacion{milcVerde}{Estación 1 de 3 · Escucha}

\begin{softbox}{milcVerde}{milcGris}{Escena para escuchar}
\textbf{Actividad 1 · IDENTIFICA --- Lo que se dice de los cuerpos.} \textbf{Sin nombres: usa iniciales.} \textbf{Primera parte.} Escribe \textbf{los comentarios sobre el cuerpo que circulan en tu salón}: apodos, chistes, «consejos» no pedidos, comentarios sobre lo que alguien come. \emph{Incluye los que se dicen «de cariño»: cuentan igual.} \textbf{Segunda parte.} ¿Sobre \textbf{qué partes} son casi siempre? ¿Le tocan igual a los hombres y a las mujeres? \emph{Anota la diferencia si la hay.} \textbf{Tercera parte, la que solo lees tú.} ¿\textbf{Qué te dices a ti mismo} sobre tu cuerpo cuando te ves en una foto o en el espejo? \emph{Escríbelo textual.} \textbf{Y después una sola pregunta:} \emph{¿le dirías eso, con esas mismas palabras, a un amigo?} \textbf{Y la última:} ¿alguna vez has hecho tú un comentario sobre el cuerpo de alguien? \emph{Cuenta cuántas veces esta semana.}
\end{softbox}

{\sffamily\bfseries\fontsize{8}{10}\selectfont\textcolor{milcNegro!55}{MARCA CUANDO LO HAYAS HECHO}}
\milccheckrow{Listaste los comentarios que circulan, incluidos los «de cariño».}
\milccheckrow{Anotaste si les tocan distinto a hombres y a mujeres.}
\milccheckrow{Escribiste textual lo que te dices a ti mismo, y si se lo dirías a un amigo.}

\begin{infoband}{milcVerde}


\textbf{En tu cuaderno (Actividad 1 · IDENTIFICA).} Título: «Actividad 1. Lo que se dice de los cuerpos». \textbf{Formato:} los comentarios que circulan + sobre qué partes + lo que te dices tú + cuántos comentarios hiciste esta semana. \textbf{Extensión:} media página. \textbf{Sabes que terminaste cuando} respondiste si le dirías a un amigo lo que te dices a ti. Esta página es tuya: \textbf{nadie más tiene que leerla si tú no quieres}.
\end{infoband}

\puente{Ya lo tienes. Ahora, quién fabricó la vara, por qué comparse es normal y qué le hicieron las pantallas a ese mecanismo.}

\milcestacion{milcTurquesa}{Fase 2 · Sistematización: entender la vara}

\begin{softbox}{milcTurquesa}{milcGris}{Cuatro claves sobre el cuerpo comparado}
Cuatro claves para dejar de discutir qué cuerpo está bien y empezar a mirar quién puso la vara.
\end{softbox}

\milcpaso{1}{milcTurquesa}{\textbf{La vara se fabricó, y se puede fechar.} El \textbf{90-60-90} no es una ley de la naturaleza: son tres números que circularon en \textbf{Cromos} alrededor del Reinado Nacional de Belleza \textbf{entre 1957 y 1962}, y que los estudios describen como la construcción de un \emph{mito}. \emph{Antes de eso, nadie en Colombia se medía así.} \textbf{Y los mismos estudios muestran de dónde salió el empuje:} el crecimiento de las industrias \textbf{textil, cosmética y de la higiene} favoreció una idea de belleza \textbf{que se pudiera comprar}. \emph{Ponlo junto con el momento 6 y el cuadro queda completo:} \textbf{primero se fabrica una vara, después se vende lo que hace falta para alcanzarla}. \textbf{La prueba está en tu propia casa:} \emph{mira fotos de hace cuarenta años ---la vara era otra, y quienes la cumplían eran otras personas---.}}
\milcpaso{2}{milcTurquesa}{\textbf{Comparse no es un defecto: es cómo funciona la cabeza.} Festinger lo planteó así: hay un \textbf{impulso a evaluar las propias capacidades}, y cuando \textbf{no existe una medida objetiva}, la evaluación se hace \textbf{comparándose con otros} (Festinger, 1954). \emph{Para correr hay cronómetro; para «estar bien» no hay ninguno.} \textbf{Por eso la instrucción «no te compares» no funciona:} es como decirle a alguien que no tenga sed. \emph{Lo que sí se puede cambiar no es \textbf{si} comparas, sino \textbf{con qué}} ---y ahí sí hay decisiones---.}
\milcpaso{3}{milcTurquesa}{\textbf{Teníamos un freno, y las pantallas se lo llevaron.} Festinger observó que \textbf{la comparación disminuye cuando la diferencia con el otro crece}: uno no se mide con quien está lejísimos. \emph{En un pueblo eso funcionaba solo ---uno se comparaba con los del barrio---.} \textbf{Un feed rompe el freno de dos maneras.} \textbf{Primera:} te pone al lado, cada día, a personas de todo el mundo escogidas por ser excepcionales. \textbf{Segunda, y peor:} \emph{esas imágenes están seleccionadas y arregladas}. \textbf{Te estás comparando con algo que no existe} ---ni siquiera para la persona de la foto---. \emph{Y ojo con quién tiene esto peor:} \textbf{el algoritmo te muestra más de lo que miras}, así que \emph{al que le duele es al que más le muestran}.}
\milcpaso{4}{milcTurquesa}{\textbf{El comentario sobre el cuerpo se queda mucho después.} La mayoría de los comentarios sobre cuerpos en un salón se dicen \textbf{sin intención de hacer daño} ---«de cariño», «es un chiste», «yo se lo digo por su bien»---. \emph{Y aun así son los que más duran:} muchas personas adultas pueden repetir, palabra por palabra, algo que les dijeron a los trece años. \textbf{Acuérdate del momento 4 del grado 7: la palabra que hiere no necesita mala intención para herir.} \emph{Y hay una categoría especialmente dañina y muy común:} \textbf{el comentario «positivo» sobre el cuerpo de otro} ---elogiar que alguien adelgazó, por ejemplo---, \emph{porque enseña a todo el salón que el cuerpo es algo que se califica}.}

\begin{drawbox}{milcTurquesa}{Las tres reglas sobre el cuerpo, propio y ajeno}
\begin{pasoslista}
\item \textbf{No comento cuerpos.} \emph{Ni para mal ni para bien: comentar «qué flaco estás» también enseña que el cuerpo se califica.}
\item \textbf{No me digo lo que no le diría a un amigo.} \emph{Si no se lo dirías a alguien que quieres, tampoco es cierto sobre ti.}
\item \textbf{Reviso con qué me estoy comparando.} \emph{Si es una foto escogida entre cien y editada, no es una comparación: es una trampa.}
\item \textbf{Y una cuarta que vale por todas:} \emph{lo que el cuerpo \textbf{hace} ---camina, juega, abraza, aguanta--- no aparece en ninguna vara.}
\end{pasoslista}
\end{drawbox}

\begin{softbox}{milcMostaza}{milcGris}{Errores comunes y su corrección}
\textbf{«No te compares»:} imposible; lo que se elige es con qué. \textbf{El elogio del cuerpo:} felicitar por adelgazar enseña que el cuerpo se califica. \textbf{«Es de cariño»:} los apodos sobre el cuerpo duran décadas. \textbf{«Se lo digo por su salud»:} nadie pidió diagnóstico. \textbf{Creer que solo les pasa a las mujeres:} la vara es distinta pero también aprieta a los hombres. \textbf{Compararse con fotos:} están escogidas entre cien y arregladas. \textbf{Pensar que la vara es natural:} tiene autor, fecha y quién gana con ella.
\end{softbox}

\begin{infoband}{milcTurquesa}


\textbf{En tu cuaderno (Actividad 2 · EXPLICA).} Título: «Actividad 2. La vara tiene fecha». \textbf{Formato:} explica con tus palabras la segunda hipótesis de Festinger y por qué una pantalla le quita el freno; después escribe \textbf{qué industrias ganan} cuando se fabrica una vara. \textbf{Extensión:} media página. \textbf{Sabes que terminaste cuando} puedes explicar por qué decirle a alguien \textbf{«no te compares» no sirve de nada}.
\end{infoband}

\puente{Ya sabes de dónde viene. Es hora de hacer dos cosas concretas: \textbf{ver de dónde salen tus comparaciones} y cambiar lo que sí puedes cambiar.}

\milcestacion{milcMagenta}{Fase 3 · Praxis: cambiar la vara}

\begin{softbox}{milcMagenta}{milcGris}{Cómo se le baja el poder a una vara}
Dos frentes: uno hacia afuera ---cómo hablamos de los cuerpos--- y otro hacia adentro ---con qué te estás midiendo---. El segundo es más difícil y el primero cambia el salón.
\end{softbox}

\milcpaso{1}{milcMagenta}{\textbf{La línea del tiempo (10 min, en grupos).} Armen una línea con \textbf{tres épocas} ---por ejemplo los años veinte, los años sesenta y hoy--- y en cada una escriban \textbf{cuál era el cuerpo considerado correcto} y \textbf{quién lo decidía} ---una revista, un reinado, una industria, un algoritmo---. \emph{Si pueden, consíganse fotos de la casa: funciona mejor que cualquier explicación.} \textbf{La conclusión sale sola:} \emph{ninguna vara duró, y todas se sintieron definitivas mientras estuvieron.}}
\milcpaso{2}{milcMagenta}{\textbf{El diario de comparaciones (una semana, 2 min al día).} Cada día anota \textbf{una vez en que te comparaste} y, al lado, \textbf{de dónde salió} la comparación: una foto en redes, alguien del salón, un comentario que te hicieron, un espejo. \emph{Al cabo de la semana, cuenta.} \textbf{Casi siempre pasa lo mismo y es información valiosísima:} la mayoría no viene de personas de tu vida real, \textbf{viene de la pantalla}.}
\milcpaso{3}{milcMagenta}{\textbf{La limpieza del feed (8 min).} Esto es lo único de esta guía que cambia algo mañana. \textbf{Abre tus redes y silencia o deja de seguir tres cuentas} que te hagan sentir mal con tu cuerpo. \emph{No hace falta pelear con nadie ni anunciarlo:} silenciar es invisible. \textbf{Y sigue dos que te muestren gente haciendo cosas} ---deporte, música, oficios, dibujo--- \textbf{en vez de gente posando.} \emph{Al algoritmo se le educa con lo que uno mira.}}
\milcpaso{4}{milcMagenta}{\textbf{Las tres reglas (7 min).} Escríbelas con tus palabras y déjalas donde las veas. \textbf{Y agrega el compromiso más difícil:} \emph{una semana entera sin comentar el cuerpo de nadie, ni para mal ni para bien.} \textbf{Lleva la cuenta de cuántas veces estuviste a punto} ---ese número, y no las reglas, es lo que te va a mostrar el tamaño de la costumbre---.}

\begin{drawbox}{milcMagenta}{Antes de cerrar tu línea de la vara}
\begin{checklista}
\item Tu línea de tiempo tiene tres épocas, con \textbf{quién decidía} en cada una.
\item Llevaste el diario de comparaciones \textbf{siete días} y contaste de dónde salieron.
\item Silenciaste o dejaste de seguir \textbf{tres cuentas} y seguiste dos de gente haciendo cosas.
\item Tus tres reglas están escritas con tus palabras y donde las veas.
\item Anotaste cuántas veces estuviste a punto de comentar el cuerpo de alguien.
\item Entendiste que \textbf{el elogio del cuerpo también califica}, y por eso también entra en la regla.
\end{checklista}
\end{drawbox}

\begin{softbox}{milcMagenta}{milcGris}{Plantilla de trabajo}
\textbf{Las tres reglas.} \emph{«No comento cuerpos ni para bien ni para mal · No me digo lo que no le diría a un amigo · Reviso con qué me comparo».} \\[1mm]
\textbf{El diario.} \emph{«Hoy me comparé con \_\_\_\_. Salió de \_\_\_\_.»} \\[1mm]
\textbf{Cuando alguien comente un cuerpo.} \emph{«Uy, con eso no» · «no hablemos de cuerpos».} ---corto, sin sermón, como en el momento 4 del grado 7---.
\end{softbox}

\begin{infoband}{milcMagenta}


\textbf{En tu cuaderno (Actividad 3 · CREA).} Título: «Actividad 3. Mi línea de la vara». \textbf{Formato:} la línea de tiempo de tres épocas + el diario de siete días con las fuentes + qué cuentas silenciaste y cuáles seguiste + las tres reglas + la cuenta de las veces que estuviste a punto. \textbf{Extensión:} una página. \textbf{Sabes que terminaste cuando} puedes decir \textbf{qué proporción de tus comparaciones vino de una pantalla}.
\end{infoband}

\puente{El producto tiene dos mitades: una mirada hacia afuera ---la vara--- y una hacia adentro ---lo que dices y lo que te dices---.}

\milcestacion{milcMostaza}{Producto de la sesión}
\textbf{Línea de la vara: tres épocas con su autor, un diario de comparaciones de siete días y tres reglas aplicadas}.
\begin{softbox}{milcMostaza}{milcGris}{Criterios de calidad}
(1) La línea de tiempo muestra tres varas distintas e identifica quién las decidía en cada época. (2) El diario cubre siete días y registra la fuente de cada comparación, no solo la comparación. (3) La limpieza del feed se hizo de verdad y se registra qué se silenció y qué se siguió. (4) Las tres reglas incluyen el comentario «positivo» sobre el cuerpo ajeno, no solo el negativo. (5) Se explica por qué «no te compares» no funciona y qué sí se puede decidir.
\end{softbox}

\puente{Antes de cerrar, revisa lo que casi siempre falta: las reglas sobre \textbf{cómo hablas del cuerpo de los demás}, no solo sobre cómo te miras tú.}

\milcestacion{milcVino}{Evaluación liberadora · cinco dimensiones}

\begin{softbox}{milcVino}{milcGris}{Sentido de la evaluación}
Evaluar no es solo revisar una nota. Es reconocer qué aprendí, qué cambió en mí, qué emociones manejé, cómo me relaciono con la ciudad y la comunidad, y qué le dejaré a quien viene detrás.
\end{softbox}

\textbf{Mi reflexión liberadora en cinco dimensiones.} Completa cada frase iniciada:

\small
\begin{tabularx}{\textwidth}{>{\bfseries\RaggedRight\arraybackslash}p{3.4cm}Y>{\RaggedRight\arraybackslash}p{4.6cm}}
\rowcolor{milcVino}\color{white}Dimensión & \color{white}\bfseries Mi respuesta (frame starter) & \color{white}\bfseries Algo que descubrí\\
1. Personal & Lo que más cambió en mí fue \linead{6.5cm} & \linead{4.2cm}\\
2. Emocional & La emoción más fuerte fue \linead{2.5cm} y la regulé \linead{3cm} & \linead{4.2cm}\\
3. Ciudadana & En democracia esto importa porque \linead{6cm} & \linead{4.2cm}\\
4. Local (Cartago) & En mi barrio sería útil para \linead{6.5cm} & \linead{4.2cm}\\
5. Intergeneracional & A mi abuela/abuelo le diría \linead{6.5cm} & \linead{4.2cm}\\
\end{tabularx}
\normalsize

\begin{infoband}{milcVino}
\textbf{Para la próxima sustentación.} Vuelve a esta tabla en seis meses. Verás que las respuestas cambian — no porque lo aprendido cambie, sino porque tú cambias. La evaluación liberadora es una conversación contigo a través del tiempo.
\end{infoband}


\puente{Cerramos con tres voces: la persona que no cabe en su apariencia, la disposición a que te muestren que estabas mal, y qué cambia cuando la imagen es editable.}

\milcestacion{milcGradeDark}{Triángulo de pensamiento · cierre filosófico}

\begin{softbox}{milcVino}{milcGris}{Dussel · lente del nosotros}
\textit{«El otro se revela realmente como otro\ldots como el pobre, el oprimido; el que, a la vera del camino, fuera del sistema, muestra su rostro sufriente y sin embargo desafiante: «¡Tengo hambre!, ¡tengo derecho a comer!»»}

{\footnotesize\itshape\color{milcNegro!70}--- Enrique Dussel · Filosofía de la liberación (1977)}

Dussel habla del \textbf{rostro}, y un rostro es lo contrario de una medida. \emph{Una vara ---90-60-90, un peso, una talla--- convierte a una persona en un número y decreta quién queda \textbf{fuera del sistema}.} \textbf{Y aquí «fuera del sistema» no es una metáfora:} el que no cumple la vara \emph{deja de ser mirado como alguien y pasa a ser mirado como un cuerpo que falla}. \textbf{Fíjate además en quién queda fuera con más facilidad:} \emph{casi siempre coincide con quien menos puede comprar lo que la vara exige} ---el gimnasio, la ropa, la comida, los tratamientos---. \textbf{La vara no es solo estética: también es económica.}

\textbf{En mi salón, ¿de quién se habla más por su cuerpo que por cualquier otra cosa?}
\end{softbox}
\linead{\linewidth}\\[1mm]
\linead{\linewidth}

\begin{softbox}{milcMostaza}{milcGris}{Estoicismo · Marco Aurelio · Meditaciones VI, 21 (c. 175 d.C.) · cuidado interior}
\textit{«Si alguien puede refutarme y probar de modo concluyente que pienso o actúo incorrectamente, de buen grado cambiaré de proceder. Pues persigo la verdad, que no dañó nunca a nadie; en cambio, sí se daña el que persiste en su propio engaño e ignorancia.»}

{\footnotesize\itshape\color{milcNegro!70}--- Marco Aurelio · Meditaciones VI, 21 (c. 175 d.C.)}

\textbf{«Se daña el que persiste en su propio engaño».} \emph{Y hay un engaño enorme funcionando aquí:} creer que la vara de hoy es \textbf{la} vara, cuando la historia muestra que cada época tuvo la suya y todas se sintieron definitivas. \textbf{La cita pide algo concreto:} \emph{estar dispuesto a cambiar de proceder cuando te muestran que estabas mal}. \textbf{Aplícalo en las dos direcciones.} \emph{Hacia afuera:} si alguien te dice que un comentario tuyo le dolió, \textbf{no discutas si era para tanto ---cambia de proceder---}. \emph{Hacia adentro:} lo que te dices frente al espejo también es una opinión, \textbf{y las opiniones se pueden refutar}.

\textbf{Si un amigo me dijera de sí mismo lo que yo me digo, ¿le diría que tiene razón?}
\end{softbox}
\linead{\linewidth}\\[1mm]
\linead{\linewidth}

\begin{softbox}{milcTurquesa}{milcGris}{Floridi · lente de la infoesfera}
\textit{«Las tecnologías digitales están re-ontologizando la naturaleza misma de la infoesfera.»}

{\footnotesize\itshape\color{milcNegro!70}--- Luciano Floridi · La cuarta revolución (2014; trad.)}

Aquí «re-ontologizar» se puede tocar con la mano: \textbf{una foto ya no es el registro de algo que pasó}. \emph{Es un objeto que se escoge entre cien, se recorta, se ilumina, se filtra y a veces se le cambia el cuerpo.} \textbf{Cambió \emph{lo que una foto es}, y sin embargo la seguimos mirando como si fuera una ventana.} \textbf{De ahí sale la regla más útil de toda esta guía:} \emph{cuando te compares con una imagen, pregúntate primero si esa imagen le corresponde a alguien}. \textbf{Muchas veces la respuesta es que no ---ni siquiera a la persona que aparece en ella---.}

\textbf{De las últimas diez fotos que vi de otra persona, ¿cuántas creo que salieron al primer intento?}
\end{softbox}
\linead{\linewidth}\\[1mm]
\linead{\linewidth}

\begin{infoband}{milcNegro}
\textbf{Nota para el docente.} Las tres son citas textuales y llevan su referencia debajo. Puedes remitir a la obra en clase.
\end{infoband}

\milcestacion{milcVino}{Mi autoevaluación · marca una casilla por fila}

{\small
\setlength{\extrarowheight}{2.2pt}
\begin{tabularx}{\textwidth}{>{\RaggedRight\arraybackslash\bfseries}p{3.0cm}YYY}
\rowcolor{milcVino}\color{white}\bfseries Criterio & \color{white}\bfseries Lo logré & \color{white}\bfseries En proceso & \color{white}\bfseries Necesito apoyo\\[.6mm]
Comprensión del tema & \milcopcion{Explico las ideas con mis palabras.} & \milcopcion{Reconozco las ideas, falta precisión.} & \milcopcion{Necesito repasar lo básico.}\\[.8mm]
\rowcolor{milcGris}Calidad del análisis & \milcopcion{Puedo mostrar que la vara cambia con la época y sé con quién me estoy comparando.} & \milcopcion{Entiendo que la moda del cuerpo cambia, pero sigo midiéndome con la de ahora.} & \milcopcion{Sigo pensando que hay cuerpos que están bien y cuerpos que están mal.}\\[.8mm]
Aplicación y producto & \milcopcion{Mi producto cumple los criterios.} & \milcopcion{Está armado, con ajustes pendientes.} & \milcopcion{Aún está en construcción.}\\[.8mm]
\rowcolor{milcGris}Reflexión 5 dimensiones & \milcopcion{Respondí cada una con honestidad.} & \milcopcion{Respondí algunas con detalle.} & \milcopcion{Mis respuestas son muy generales.}\\[.8mm]
Triángulo de pensamiento & \milcopcion{Conecté las 3 voces con mi trabajo.} & \milcopcion{Conecté al menos una.} & \milcopcion{No las relaciono con mi proceso.}\\[.8mm]
\end{tabularx}}

\vspace{3mm}
\textbf{Hoy entendí que:}\\[1mm]
\linead{\linewidth}\\[1mm]
\linead{\linewidth}

\textbf{Mi compromiso para la próxima sesión es:}\\[1mm]
\linead{\linewidth}\\[1mm]
\linead{\linewidth}

\begin{softbox}{milcMostaza}{milcGris}{Cita de despedida · Marco Aurelio}
\textit{``El obstáculo en el camino se vuelve el camino. Lo que estorba al camino se vuelve camino.''}\\[1mm]
Cada error de hoy, bien observado, se convierte en pieza del camino que pisarás mañana.
\end{softbox}

\small
\begin{tabularx}{\textwidth}{>{\bfseries}p{3.6cm}Y}
\rowcolor{milcGradeDark!12}Nombre completo: & \linead{10cm}\\
Fecha: & \linead{4cm} \quad Curso/grupo: \linead{4cm}\\
Mi firma: & \linead{9cm}\\
\end{tabularx}
\normalsize

\begin{infoband}{milcGradeDark}
\textbf{Para la próxima sesión.} Dos cosas. \textbf{Primero, cumple la semana sin comentar el cuerpo de nadie} ---ni para mal ni para bien--- y trae \textbf{el número de veces que estuviste a punto}. \emph{Trae también el diario de comparaciones, con la cuenta de cuántas vinieron de una pantalla.} \textbf{Segundo, pregunta en tu casa por la vara de antes:} averigua con alguien mayor \textbf{cómo era el cuerpo que se consideraba bonito cuando tenía tu edad}, \textbf{dónde lo veía} ---revistas, novelas, reinados, el cine--- y si alguien de su familia sufrió por no parecerse a eso. \emph{Si pueden, miren juntos una foto de esa época.} \textbf{Trae:} tus dos números y lo que te contaron. \emph{Vas a encontrar la prueba que ninguna explicación reemplaza: la vara de esa foto ya no la usa nadie, y sin embargo alguien lloró por ella.}
\end{infoband}

\end{document}
